\documentclass[10pt]{article}
\usepackage[letterpaper,margin=0.60in]{geometry}
\usepackage[T1]{fontenc}
\usepackage{lmodern}
\usepackage{microtype}
\usepackage{amsmath,amssymb}
\usepackage{enumitem}
\usepackage{needspace}
\usepackage{titlesec}
\usepackage{xcolor}
\usepackage{hyperref}

\hypersetup{colorlinks=true,linkcolor=black,urlcolor=blue,citecolor=black,
  pdftitle={Limits on Understanding That Emerge from Scale},pdfauthor={J. R. Landers}}
\setlength{\parindent}{0pt}
\setlength{\parskip}{0.32em}
\setlist{nosep,leftmargin=1.3em}
\titleformat{\section}{\Large\bfseries}{}{0pt}{}
  [\vspace{0.55em}{\color{black}\titlerule}]
\newcommand{\E}{\mathbb E}
\newcommand{\Prob}{\mathbb P}
\newcommand{\avgint}{\mathop{\rlap{\raisebox{.15ex}{--}}\!\int}}

\title{\vspace{-2.3em}\textbf{Limits on Understanding That Emerge from Scale}}
\author{J. R. Landers}
\date{September 8, 2026}

\begin{document}
\maketitle
\vspace{-1.5em}
\small
\setlength{\abovedisplayskip}{6pt plus 2pt minus 1pt}
\setlength{\belowdisplayskip}{6pt plus 2pt minus 1pt}
\setlength{\abovedisplayshortskip}{4pt plus 1pt}
\setlength{\belowdisplayshortskip}{5pt plus 1pt minus 1pt}

\textbf{Question.} Suppose an experiment can observe physical structure only up
to a length $R$. Should we expect related structure at lengths larger than $R$?
This note separates two ways to answer. A probability model can assign a chance
to larger scales. A fluid equation can explain how energy reaches them. Neither
approach allows a finite observation, by itself, to prove continuation to
arbitrarily large scales.

\textbf{Running example.} To keep the discussion concrete, consider a
two-dimensional flow forced at length $L_f=1\,\mathrm m$. The forcing supplies
energy at rate $\varepsilon=10^{-3}\,\mathrm{m^2\,s^{-3}}$, and our experiment
can detect coherent structures only up to $R=8\,\mathrm m$. We will ask what,
if anything, can be said about structures larger than $8\,\mathrm m$.

\section*{1. A probability law for further scales}
\subsection*{Scales as a point process}
Suppose that every structure we care about, such as an eddy, wave, or
instability, has a characteristic length $L$. We record only that length. The
collection of structures is then represented by a random set of points on the
positive $L$-axis, starting at a chosen minimum length $L_{\min}$.

It is natural to measure this axis logarithmically: going from 1 to 2 metres is
the same relative change as going from 100 to 200 metres. We now make a specific
modeling assumption. Counts in nonoverlapping scale intervals are independent,
and the expected number in a short interval near $L$ is
\[
  d\mu(L)=\eta(L)\frac{dL}{L},\qquad \eta(L)\ge0.
\]
These two statements define a Poisson point process. The function $\eta(L)$ is
the expected density of structures per unit of log scale. A structure is
``relevant'' only if it interacts with the phenomenon being studied. If one
instead measures relative $d$-dimensional volume, the extra factor $d$ can be
absorbed into $\eta$ and does not change the argument.

The number $N(R,S)$ of processes with scales in $(R,S]$ is Poisson with mean
\[
  \mu(R,S)=\int_R^S\eta(L)\frac{dL}{L}.
\]
The chance of seeing at least one such structure is therefore
\[
 \Prob\{N(R,S)\ge1\}
 =1-\exp\!\left[-\int_R^S\eta(L)\frac{dL}{L}\right].                   \tag{1}
\]
Now let the upper limit $S$ grow without bound. The model gives the exact test
\[
 \boxed{\text{unbounded process scales almost surely}}
 \quad\Longleftrightarrow\quad
 \boxed{\displaystyle\int_R^\infty\eta(L)\frac{dL}{L}=\infty}.       \tag{2}
\]

Here ``almost surely'' means with probability one under the model.

\subsection*{Coin-flip intuition}
First suppose $\eta(L)=\eta_0>0$ at every scale.
Each doubling interval
$(2^jR,2^{j+1}R]$ has the same hit probability
\[
 p=1-\exp(-\eta_0\log2)=1-2^{-\eta_0}.
\]
The intervals are independent, so the probability of missing the first $m$ is
$2^{-m\eta_0}\to0$. This is like repeating a coin toss with a fixed positive
chance of heads: eventually a head appears, and in fact heads appear infinitely
often. Likewise, structures occur in infinitely many doubling intervals with
probability one. Arratia calls this a scale-invariant Poisson process
\cite{arratia}.

In the running example, take $\eta_0=1$. Each interval
$(8,16],(16,32],(32,64]$, and $(64,128]$ metres then has probability $1/2$ of
containing a relevant structure. The probability of finding at least one by
$128\,\mathrm m$ is
\[
  1-\left(\frac12\right)^4=0.9375.
\]
If the same constant intensity continues through infinitely many doubling
intervals, structures occur at arbitrarily large scales with probability one.

\subsection*{A decaying tail}
The answer changes if large structures become sufficiently rare. For example,
if
$\eta(L)=\eta_0(L_{\min}/L)^a$ with $a>0$, then
\[
 \int_R^\infty\eta(L)\frac{dL}{L}
 =\frac{\eta_0}{a}\left(\frac{L_{\min}}R\right)^a<\infty.
\]
then the total expected number beyond $R$ is finite. The probability that no
relevant structure exists beyond $R$ is therefore
\[
 \exp\!\left[-\frac{\eta_0}{a}
 \left(\frac{L_{\min}}R\right)^a\right]
\]

For a numerical contrast, define two models that agree throughout the observed
range:
\[
 \eta_A(L)=1,\qquad
 \eta_B(L)=
 \begin{cases}
  1, & L\le8\,\mathrm m,\\
  (8\,\mathrm m/L)^2, & L>8\,\mathrm m.
 \end{cases}
\]
Model $A$ predicts unbounded scales almost surely. For model $B$,
\[
 \int_{8\,\mathrm m}^{\infty}\left(\frac{8\,\mathrm m}{L}\right)^2
 \frac{dL}{L}=\frac12,
 \qquad
 \Prob\{\text{no structure beyond }8\,\mathrm m\}=e^{-1/2}\approx0.607.
\]

\subsection*{Why finite data cannot decide}
Models $A$ and $B$ fit every possible
observation made at $L\le8\,\mathrm m$, yet they give opposite answers about
unlimited continuation. The answer is determined by the assumed tail of
$\eta$, which the finite experiment cannot see.

\section*{2. What Navier--Stokes can supply}
The Poisson model describes how many scales occur, but it does not explain what
creates them. Two-dimensional fluid flow provides a possible mechanism: an
inverse cascade can carry energy from the forcing scale toward larger lengths.
Consider the randomly forced, damped Navier--Stokes equations on a periodic
square $\mathbb T_\lambda^2$ of size $\lambda$:
\[
 du+\big[(u\!\cdot\!\nabla)u+\nabla p\big]dt
 =\big[\nu\Delta u-\alpha(-\Delta)^{-2\gamma}u\big]dt+dW^\lambda,
 \qquad \nabla\!\cdot u=0.                                             \tag{3}
\]
Here $u$ is velocity, $p$ is pressure, $\nu$ is viscosity, $\alpha$ is
large-scale drag, and $W^\lambda$ is the random forcing. The forcing injects
mean energy $\varepsilon>0$ per unit time and area near one fixed length.

\subsection*{The inverse-flux law}
Bedrossian, Coti Zelati, Punshon-Smith, and Weber proved an inverse-flux law for
stationary solutions as the box grows and viscosity and drag vanish
\cite{bcpw}. Among their assumptions are
\[
\nu\E\|\omega\|_2^2\to0,
 \qquad
 \alpha\E\|(-\Delta)^{-\gamma}\omega\|_2^2\to0.                    \tag{4}
\]
where $\omega=\operatorname{curl}u$ is vorticity. The powers $-2\gamma$ and
$-\gamma$ are consistent: a quadratic energy norm splits the operator equally
between its two factors, as in Ref.~\cite{bcpw}.

To measure transfer across a length $\ell$, compare velocities at two points
separated by $\ell$ in direction $n$:
$\delta_{\ell n}u(x)=u(x+\ell n)-u(x)$. Average the resulting cubic quantity
over all positions $x$ and directions $n$:
\[
 \mathcal F_\ell(u)=
 \avgint_{S^1}\avgint_{\mathbb T_\lambda^2}
 |\delta_{\ell n}u|^2(\delta_{\ell n}u\cdot n)\,dx\,dn.
\]
The theorem gives an upper inertial-range length $\ell_\alpha\to\infty$ for
which
\[
 \E\mathcal F_\ell(u)\sim2\varepsilon\ell,
 \qquad 1\ll\ell\ll\ell_\alpha.                                    \tag{5}
\]
The positive sign says that energy is moving toward larger lengths. Thus, for
any chosen finite $R$, one can take a large enough box and weak enough drag to
obtain active inverse transfer at some $\ell\ge R$. Notice the quantifiers: the
system may change when $R$ changes. This does not yet describe one flow whose
active scale grows forever.

\Needspace{10\baselineskip}
\subsection*{Numerical flux at 8 metres}
For the running example, suppose $8\,\mathrm m$ lies inside the inertial range.
Equation (5) then gives
\[
 \E\mathcal F_{8\,\mathrm m}
 \approx2(10^{-3}\,\mathrm{m^2\,s^{-3}})(8\,\mathrm m)
 =0.016\,\mathrm{m^3\,s^{-3}}.
\]
This positive number says that energy is crossing the $8\,\mathrm m$ scale in
the direction of still larger structures. It supplies a physical mechanism
for continuation beyond the measurement limit. It does not prove that the
scale of one particular flow grows without bound.

\subsection*{The role of ergodicity}
There is also an almost-sure statement at each fixed set of parameters. If the
chosen stationary measure is ergodic, Birkhoff's theorem gives
\[
 \frac1T\int_0^T\mathcal F_\ell(u(t))\,dt
 \longrightarrow \E\mathcal F_\ell(u)>0.                              \tag{6}
\]
So the long-time average flux is positive for almost every realization. This
is a stationary result for an allowed scale in a fixed box. It does not remove
the box's largest possible length.

\section*{3. The stronger theorem still wanted}
The stronger goal concerns one flow on the whole plane $\mathbb R^2$. A
spatially homogeneous flow on an infinite plane generally has infinite total
energy, so the right object is a homogeneous statistical solution with finite
energy \emph{per unit area}, not a finite-energy $L^2(\mathbb R^2)$ solution.

\subsection*{Defining the large scale}
We must also define what the flow's ``large scale'' means. Let $E(k,t)$ be its
energy spectrum, where small wave number $k$ means large physical length. Fix a
fraction $\theta\in(0,1)$ and define
\[
 k_\theta(t)=\inf\!\left\{K>0:\int_0^K E(k,t)\,dk\ge
 \theta\int_0^\infty E(k,t)\,dk\right\},\qquad L(t)=k_\theta(t)^{-1}.
\]
Thus $k_\theta(t)$ is the wave number below which the chosen fraction of energy
lies, and $L(t)=1/k_\theta(t)$ is its associated length. Energy moving toward
smaller $k$ makes $L(t)$ larger. The desired theorem is
\[
 \Prob\{L(t)\to\infty\}=1.                                             \tag{7}
\]
\subsection*{A conditional growth law}
What growth rate might we expect? If we additionally assume Kraichnan's
inverse-cascade spectrum~\cite{kraichnan}, with forcing length $L_f$,
\[
 E(k)=C_I\varepsilon^{2/3}k^{-5/3},
 \qquad L(t)^{-1}\le k\le L_f^{-1}.
\]
then the energy accumulated by time $t$ is approximately the integral of this
spectrum. Setting that integral equal to $\varepsilon t$ gives
\[
 \varepsilon t
 =\frac{3C_I}{2}\varepsilon^{2/3}
   \big(L(t)^{2/3}-L_f^{2/3}\big),
\]
so
\[
\boxed{L(t)=\left[L_f^{2/3}+
 \frac{2}{3C_I}\varepsilon^{1/3}t\right]^{3/2}}
 \sim \left(\frac{2}{3C_I}\right)^{3/2}\varepsilon^{1/2}t^{3/2}.       \tag{8}
\]
For a simple numerical illustration, set $C_I=1$. This choice makes the
arithmetic transparent; it is not an empirical calibration. Using the values
from the running example, (8) becomes
\[
 L(t)=1\,\mathrm m\left(1+\frac{t}{15\,\mathrm s}\right)^{3/2}.
\]
The predicted scale reaches the observation limit after $45$ seconds,
\[
 L(45\,\mathrm s)=1\,\mathrm m\,(1+3)^{3/2}=8\,\mathrm m,
\]
and reaches $27\,\mathrm m$ after $120$ seconds. These numbers illustrate what
the $t^{3/2}$ law says; they are conditional on the assumed spectrum.

\Needspace{9\baselineskip}
\subsection*{What remains open}
This calculation predicts $L(t)\propto t^{3/2}$, but it is not a theorem derived
from (3): it assumes the spectrum that produces the answer. The lower spectral
edge used in the calculation is proportional to the quantile length above, up
to a constant that does not change the exponent.

Finally, this growing state is not stationary, so the Birkhoff argument in (6)
cannot prove (7) or (8). Drag, finite forcing time, or a container may stop the
growth. The open problem is to prove from Navier--Stokes that the infrared scale
of one unbounded flow grows almost surely, under explicit assumptions, without
first assuming the $k^{-5/3}$ law.

\begin{thebibliography}{9}\footnotesize
\setlength{\itemsep}{1pt}
\bibitem{arratia} R.~Arratia, ``On the Central Role of the Scale Invariant
Poisson Processes on $(0,\infty)$,'' DIMACS Series \textbf{41} (1998),
\href{https://arxiv.org/abs/1611.05572}{arXiv:1611.05572}.
\bibitem{bcpw} J.~Bedrossian, M.~Coti Zelati, S.~Punshon-Smith, and F.~Weber,
``Sufficient Conditions for Dual Cascade Flux Laws in the Stochastic 2d
Navier--Stokes Equations,'' \emph{Arch. Ration. Mech. Anal.} \textbf{237}
(2020), \href{https://doi.org/10.1007/s00205-020-01503-9}{103--145}.
\bibitem{kraichnan} R.~H.~Kraichnan, ``Inertial Ranges in Two-Dimensional
Turbulence,'' \emph{Phys. Fluids} \textbf{10} (1967),
\href{https://doi.org/10.1063/1.1762301}{1417--1423}.
\end{thebibliography}

\end{document}
